\section{Deep Q-Learning and Model-Based RL}

\subsection{DQN: Scaling up with Neural Networks}
The \textbf{Deep Q-Network (DQN)} was a landmark algorithm that successfully combined Q-learning with Deep Neural Networks. To overcome the instabilities of the "Deadly Triad," DQN introduced two crucial mechanisms:
\begin{enumerate}
    \item \textbf{Experience Replay:} Instead of updating the network with the latest transition, the agent stores its experiences in a buffer and samples mini-batches uniformly. This breaks the temporal correlation between successive steps and allows for data reuse.
    \item \textbf{Target Networks:} DQN uses two networks. The "online" network is updated constantly, while a "target" network with frozen weights is used to compute the Q-learning targets. This prevents the target from "moving" too quickly, which would otherwise cause oscillations or divergence.
\end{enumerate}

\subsection{Evolving the DQN Architecture}
Since the original DQN, several improvements have become standard:
\begin{itemize}
    \item \textbf{Double DQN:} Addresses the \textit{Maximization Bias}. It uses the online network to select the best action and the target network to evaluate its value.
    \item \textbf{Dueling DQN:} Instead of estimating $Q(s, a)$ directly, the network architecture is split into two streams: one for the state-value $V(s)$ and one for the action advantage $A(s, a)$. This allows the agent to learn which states are valuable without having to learn the effect of every action in that state.
    \item \textbf{Prioritized Experience Replay (PER):} Instead of sampling uniformly, PER samples transitions that have a high TD-error, assuming they contain more "surprising" or useful information for learning.
\end{itemize}

\subsection{Model-Based RL: Planning and Learning}
While Model-Free RL learns from raw experience, \textbf{Model-Based RL} attempts to learn the environment's dynamics ($P$ and $R$). If the agent has a model, it can "think" before it acts—a process we call \textbf{Planning}.
The \textbf{Dyna-Q} architecture provides a beautiful integration of both. The agent uses real experience to update its value function (\textit{Learning}) and to improve its model. Simultaneously, it uses its model to generate simulated experiences (\textit{Planning}) to further refine its value function.

\subsection{Simulation-Based Search and MCTS}
For extremely complex games (like Go), we use \textbf{Monte Carlo Tree Search (MCTS)}. Unlike global planning, MCTS performs \textit{local} planning by building a search tree from the current state. It uses four steps: selection (using a tree policy like UCT), expansion, simulation (rollout), and backup. This allowed AlphaGo to evaluate moves without needing to explore the entire state space.

In continuous environments, we encounter the problem of infinite branching. Techniques like \textbf{Progressive Widening} are used here to dynamically control how many new actions or next-states are added to the search tree as more simulations are performed.
 Greenland
